%% =============================================================================
%% chapters/01-foundations.tex
%% PART I: Foundations and Concepts
%% Chapter 1: Platform Motivation, HIE Context & Architectural Principles
%% =============================================================================

\chapter{Platform Motivation, HIE Context \& Architectural Principles}
\label{chap:foundations}

The Motivation and Strategy layer articulates the fundamental healthcare interoperability drivers, stakeholder concerns, strategic goals, architecture principles, and business capabilities that shape the Harmonia Health Information Exchange (HIE) platform.

\begin{figure}[htbp]
    \centering
    \begin{adjustbox}{max width=\textwidth}
        %% =============================================================================
%% diagrams/fig-motivation-map.tex
%% ArchiMate Motivation & Strategy Architecture View
%% =============================================================================
\begin{tikzpicture}[node distance=1.0cm and 0.8cm, auto]

    % --- Background Layer Plates ---
    \draw[archi-plate-bg] (-0.8, 7.5) rectangle (15.8, 9.7);
    \node[archi-plate-title] at (-0.6, 9.5) {Stakeholder Layer};

    \draw[archi-plate-bg] (-0.8, 4.0) rectangle (15.8, 7.2);
    \node[archi-plate-title] at (-0.6, 7.0) {Drivers \& Strategic Goals};

    \draw[archi-plate-bg] (-0.8, 2.2) rectangle (15.8, 3.8);
    \node[archi-plate-title] at (-0.6, 3.6) {Principles \& Requirements};

    \draw[archi-plate-bg] (-0.8, 0.4) rectangle (15.8, 2.0);
    \node[archi-plate-title] at (-0.6, 1.8) {Core Capabilities};

    % --- 1. Stakeholders ---
    \node[archi-stakeholder, minimum width=3.4cm] (stk_hie)  at (0.8, 8.4) {\archibadge{Stakeholder}\archiname{Health Network}\\\archidesc{Regional Operator}};
    \node[archi-stakeholder, minimum width=3.4cm] (stk_prov) at (5.0, 8.4) {\archibadge{Stakeholder}\archiname{Healthcare Provider}\\\archidesc{Hospitals \& Clinics}};
    \node[archi-stakeholder, minimum width=3.4cm] (stk_clin) at (9.2, 8.4) {\archibadge{Stakeholder}\archiname{Care Team}\\\archidesc{Clinicians \& Nurses}};
    \node[archi-stakeholder, minimum width=3.4cm] (stk_pat)  at (13.4, 8.4) {\archibadge{Stakeholder}\archiname{Patient}\\\archidesc{Care Recipient}};

    % --- 2. Drivers ---
    \node[archi-driver, minimum width=4.5cm] (drv_interop) at (1.8, 6.0) {\archibadge{Driver}\archiname{Clinical Interoperability}\\\archidesc{HL7 v2 to FHIR R5 Unified Model}};
    \node[archi-driver, minimum width=4.5cm] (drv_uptime)  at (7.5, 6.0) {\archibadge{Driver}\archiname{24/7 Availability}\\\archidesc{Zero Downtime Fault Tolerance}};
    \node[archi-driver, minimum width=4.5cm] (drv_privacy) at (13.2, 6.0) {\archibadge{Driver}\archiname{Audit \& PHI Privacy}\\\archidesc{Strict Security \& Compliance}};

    % --- 3. Goals ---
    \node[archi-goal, minimum width=4.5cm] (gol_realtime) at (1.8, 4.6) {\archibadge{Goal}\archiname{Real-Time Ingestion}\\\archidesc{Sub-second Pipeline Ingestion}};
    \node[archi-goal, minimum width=4.5cm] (gol_zeroloss) at (7.5, 4.6) {\archibadge{Goal}\archiname{Zero Message Loss}\\\archidesc{Durable Clustered Transport}};
    \node[archi-goal, minimum width=4.5cm] (gol_subms)    at (13.2, 4.6) {\archibadge{Goal}\archiname{Sub-ms Cache Access}\\\archidesc{Distributed Low-Latency Grid}};

    % --- 4. Principles & Requirements ---
    \node[archi-principle, minimum width=3.5cm]   (prn_atleastonce) at (0.9, 3.0) {\archibadge{Principle}\archiname{At-Least-Once}\\\archidesc{Idempotent Processing}};
    \node[archi-requirement, minimum width=3.5cm] (req_fhir5)       at (5.1, 3.0) {\archibadge{Requirement}\archiname{FHIR R5 Schemas}\\\archidesc{14 Core Resource Models}};
    \node[archi-requirement, minimum width=3.5cm] (req_ha)          at (9.3, 3.0) {\archibadge{Requirement}\archiname{Artemis HA Failover}\\\archidesc{Automatic <2s Recovery}};
    \node[archi-principle, minimum width=3.5cm]   (prn_writebehind) at (13.5, 3.0) {\archibadge{Principle}\archiname{Write-Behind}\\\archidesc{Decoupled Relational IO}};

    % --- 5. Capabilities ---
    \node[archi-capability, minimum width=3.5cm] (cap_ingest)   at (0.9, 1.2) {\archibadge{Capability}\archiname{Gateway Ingestion}\\\archidesc{Pylai MLLP / REST}};
    \node[archi-capability, minimum width=3.5cm] (cap_msg)      at (5.1, 1.2) {\archibadge{Capability}\archiname{Clustered Messaging}\\\archidesc{Petasos HA Backbone}};
    \node[archi-capability, minimum width=3.5cm] (cap_workflow) at (9.3, 1.2) {\archibadge{Capability}\archiname{Modular Workflows}\\\archidesc{Ponos / Praxis / Erga}};
    \node[archi-capability, minimum width=3.5cm] (cap_grid)     at (13.5, 1.2) {\archibadge{Capability}\archiname{In-Memory Cache}\\\archidesc{Hestia Mneme Grid}};

    % --- Connecting Relationships ---
    % Stakeholder -> Driver (Association)
    \draw[archi-association] (stk_hie) -- (drv_interop);
    \draw[archi-association] (stk_prov) -- (drv_uptime);
    \draw[archi-association] (stk_clin) -- (drv_interop);
    \draw[archi-association] (stk_pat) -- (drv_privacy);

    % Driver -> Goal (Influence)
    \draw[archi-influence] (drv_interop) -- (gol_realtime);
    \draw[archi-influence] (drv_uptime) -- (gol_zeroloss);
    \draw[archi-influence] (drv_privacy) -- (gol_subms);

    % Goal -> Principle / Requirement (Realization)
    \draw[archi-realization] (prn_atleastonce) -- (gol_zeroloss);
    \draw[archi-realization] (req_fhir5) -- (gol_realtime);
    \draw[archi-realization] (req_ha) -- (gol_zeroloss);
    \draw[archi-realization] (prn_writebehind) -- (gol_subms);

    % Principle/Requirement -> Capability (Realization)
    \draw[archi-realization] (cap_ingest) -- (req_fhir5);
    \draw[archi-realization] (cap_msg) -- (prn_atleastonce);
    \draw[archi-realization] (cap_msg) -- (req_ha);
    \draw[archi-realization] (cap_workflow) -- (req_fhir5);
    \draw[archi-realization] (cap_grid) -- (prn_writebehind);

\end{tikzpicture}

    \end{adjustbox}
    \caption{ArchiMate Motivation and Strategy Architecture View for Harmonia HIE}
    \label{fig:foundations_motivation_map}
\end{figure}

\section{Healthcare Interoperability Context \& Stakeholders}

Modern regional healthcare networks operate in heterogeneous environments characterized by disparate legacy protocols (HL7 v2.3--v2.5 over MLLP/TCP), emergent web standards (HL7 FHIR R5 over HTTPS/REST), and distributed operational requirements. Harmonia provides an enterprise-grade integration bus and canonical state engine engineered to serve key stakeholders:

\begin{description}
    \item[\textbf{Regional Health Network Operator}] Responsible for governing regional health data exchange, ensuring service-level agreement (SLA) adherence, managing cross-facility data governance, and maintaining uninterrupted interoperability bridges.
    \item[\textbf{Healthcare Providers (Hospitals, Health Systems, Diagnostic Labs)}] Organizations that produce and consume high volumes of clinical messages (e.g. ADT admissions, ORU lab observations, MDM clinical notes) via legacy MLLP streams and modern REST APIs.
    \item[\textbf{Healthcare Directory Stewards \& Credentialing Authorities}] Governance officers, medical registrars, and directory administrators charged with validating practitioner qualifications, managing organizational hierarchies, verifying electronic service endpoints, and maintaining master practitioner registries.
    \item[\textbf{Clinicians \& Care Teams}] Physicians, nurses, and allied health staff requiring immediate, accurate, and consolidated longitudinal patient records and up-to-date provider directories at the point of care.
    \item[\textbf{Patients \& Care Recipients}] The ultimate beneficiaries of timely, coordinated healthcare delivery whose Protected Health Information (PHI) must be safeguarded with absolute confidentiality.
\end{description}

\section{Core Architectural Drivers}

The platform architecture resolves five foundational drivers:
\begin{enumerate}
    \item \textbf{Clinical Interoperability Mandates}: Seamless translation of legacy HL7 v2 triggers into standardized HL7 FHIR Release 5 schemas without loss of clinical nuance or diagnostic context.
    \item \textbf{National Healthcare Directory Interoperability}: Maintenance of synchronized, standards-compliant FHIR R5 provider directories (\texttt{Practitioner}, \texttt{PractitionerRole}, \texttt{Organization}, \texttt{Location}, \texttt{HealthcareService}, \texttt{Endpoint}, \texttt{Group}) with strict referential integrity and national provider identifier governance (e.g., Australian HPI-I, HPI-O).
    \item \textbf{24/7/365 High-Availability Operations}: High-availability messaging clusters and cache grids ensuring zero message loss, sub-second failover, and zero pipeline stalls during emergency admissions.
    \item \textbf{Patient Safety \& Identity Integrity}: Automated, deterministic Enterprise Master Patient Index (EMPI) reconciliation preventing duplicate or mismatched patient identities.
    \item \textbf{PHI Privacy \& Regulatory Compliance}: Strict adherence to HIPAA, GDPR, and regional privacy frameworks requiring encrypted transport, audit trails (\texttt{AuditEvent}, \texttt{Provenance}), and zero PHI leakage into operational logs.
\end{enumerate}

\section{Core Architectural Principles}

\begin{principlebox}{At-Least-Once Delivery \& Idempotent Processing}
The platform assumes at-least-once message delivery across all messaging transports. Consuming components must be idempotent. Deduplication is enforced at broker level (via ActiveMQ Artemis \texttt{\_AMQ\_DUPL\_ID} caches), client transport level (\texttt{DuplicateDetector}), and business persistence level.
\end{principlebox}

\begin{principlebox}{Separation of Transport from Business Logic}
The messaging framework (\texttt{Petasos}) treats all payloads as opaque binary/text streams. Transport routes are unaware of clinical schemas. Transformation logic is strictly isolated in modular task activities (\texttt{Erga}).
\end{principlebox}

\begin{principlebox}{Governed Asynchronous Change Management for Directory Entities}
HTTP \texttt{POST} and \texttt{PUT} interactions on directory resources represent requests for change rather than immediate direct mutations of authoritative state. Every write request generates an immutable \texttt{ProviderRegistryChangePragma}, undergoes validation and automated referential verification in the \texttt{Energeia} pipeline, and commits only upon successful verification.
\end{principlebox}

\begin{principlebox}{Write-Behind In-Memory Cache Persistence}
Workflow pipelines write state snapshots exclusively to the in-memory cache grid (\texttt{Mneme}) to prevent disk I/O bottlenecks. Mneme asynchronously persists snapshots to relational storage (\texttt{Mnemosyne}) via non-blocking write-behind cache stores.
\end{principlebox}

\begin{principlebox}{Canonical Modeling with Standardized FHIR R5 Boundaries}
Internal workflow state is represented using the canonical \texttt{Pragma} domain model in \texttt{Calliope}. All external boundary interactions and REST APIs adhere strictly to HL7 FHIR R5 resource definitions.
\end{principlebox}

\begin{principlebox}{Governed PHI Logging \& Dual-Logger Privacy Controls}
Operational log streams (\texttt{INFO}, \texttt{WARN}, \texttt{ERROR}) must remain strictly devoid of Protected Health Information (PHI), personally identifiable demographics, and security credentials under all conditions. Clinical payload data and diagnostic identifiers are permitted exclusively at \texttt{DEBUG} and \texttt{TRACE} severity levels via the dedicated \texttt{PhiLogger} dual-gate security evaluation mechanism, routing non-additively to isolated, encrypted diagnostic storage.
\end{principlebox}

\section{Strategic Value Stream Flow}

The platform executes clinical event workflows across five coordinated processing phases:
\begin{equation*}
\begin{aligned}
\text{Ingest Trigger} &\longrightarrow \text{Transport \& Deduplicate} \longrightarrow \text{Orchestrate Erga} \\
                      &\longrightarrow \text{Update Cache Grid} \longrightarrow \text{Dispatch Egress / Serve FHIR APIs}
\end{aligned}
\end{equation*}

\newpage
